\documentclass[exercice,a4]{thedocument}%
\usepackage{pstricks}%
\startdatakeys
    \source{25GENMATAS1}
    \BO{2026}
    \cycle{4}
    \competencesDuSocle{CA,RA}
    \objectifsApprentissage{C4A1aa,C4A2ad,C4B2cb,C4D2aa,C4D3ac}
    \motsCles{droites parallèles,Thalès,Pythagore,aire}
\enddatakeys
\begin{document}%
\providecommand{\degres}{^\circ}%
    

On considère la figure ci-dessous qui n'est pas représentée en vraie grandeur.

\medskip

\begin{minipage}{.54\textwidth}
On donne les informations suivantes :

\begin{itemize}
\item (BD) et (AC) sont perpendiculaires.
\item (AD) et (AB) sont perpendiculaires.
\item (AB) et (BC) sont perpendiculaires.
\item AE $= 9,6$\,cm ; CE $= 5,4$\,cm ; 
BC $= 9$\,cm.
\end{itemize}
\end{minipage} \hfill
\begin{minipage}{.41\textwidth}
\psset{unit=.3cm,arrowsize=2pt 3}
\begin{pspicture}(-2,-2)(17,13)
%\psgrid
\pspolygon(0,12)(16,12)(0,0)
\pspolygon(0,0)(9,0)(0,12)
\psframe(0,0)(1,1)\psframe(0,12)(1,11)
\rput{35}(5.8,4.25){\psframe(1,1)}
\uput[u](0,12){A} \uput[dl](0,0){B} \uput[dr](9,0){C} \uput[ul](16,12){D} \uput[r](6,4.3){E}
\end{pspicture}
\end{minipage}

\medskip

\begin{enumerate}
\item Montrer que les droites (AD) et (BC) sont parallèles.
\item Calculer la longueur AD.
\item Montrer que la longueur BE est de $7,2$\,cm.
\item Est-il vrai que l'aire du triangle ABE représente le tiers de l'aire du triangle ABD ?
\end{enumerate}\hfill\parbox{.5\textwidth}{}
\showthesource
\end{document}
